\section{Related Work}
\label{sec:related}

\subsection{Efficient Attention}
\label{sec:related-efficient}

Reducing the quadratic cost of self-attention is a central research theme, especially in the era of long-horizon agents.
\emph{Training-free sparse} methods introduce sparsity at inference via fixed patterns, heuristic eviction strategies, or lightweight importance estimation~\citep{h2o,streamingllm,sparq,quest,minference,duoattention,xattention,flexprefill,moa,sampleattention,lserve,spargeattention}.
Nevertheless, the resulting training-inference mismatch can cause error accumulation in long-context settings~\citep{lil}.
In contrast, \emph{trainable sparse} methods incorporate sparsity directly into the training stage, for example, through learned gating mechanisms~\citep{seerattention,seerattentionr}, end-to-end sparse pre-training~\citep{nsa}, block-level mixture routing~\citep{moba,infllm-v2,minicpm4}, or full-to-sparse distillation~\citep{deepseekv32,ssa}.
DSA~\citep{deepseekv32}, the foundation of our work, distills a lightweight lightning indexer from full attention to select the top-$k$ tokens for each query, reducing the core attention complexity to~$O(Lk)$.
Beyond sparsity, \emph{hybrid architectures} reduce the number of expensive quadratic layers by interleaving them with sliding window attention (with or without sink)~\citep{gptoss,gemma3,mimov2,longcat}, linear attention~\citep{gdn,minimax01,nemotron-h,kimilinear}, or state-space layers~\citep{mamba,mamba2,jamba}.

\subsection{Cross-Layer Sharing}
\label{sec:related-crosslayer}
Recent studies demonstrate that representation exhibits strong consistency across adjacent layers.
This structural property is often leveraged to reduce computational redundancy and accelerate inference.
TidalDecode~\citep{tidaldecode}, LessIsMore~\citep{lessismore}, OmniKV~\citep{omnikv}, and DELTA~\citep{delta} reuse top-$k$ indices from periodic anchor layers for sparse decoding.
Kascade~\citep{kascade} formalizes anchor layer selection via dynamic programming over a cross-layer similarity matrix and identifies head-aware remapping as critical for maintaining accuracy.
All of these methods rely on \emph{full attention} at anchor layers to compute exact top-$k$ indices.
Independently, cross-layer \emph{KV cache sharing} reduces memory by letting multiple layers reuse the same key-value tensors~\citep{yoco,cla,minicache,swiftkv,mlkv,wu2025improving}.
HySparse~\citep{hysparse} unifies both directions, interleaving full attention layers with sparse layers that inherit both top-$k$ block indices and KV~caches.
However, all of these approaches require full attention layers as the oracle, which DSA removes completely.

IndexCache differs in two key aspects.
First, the oracle is fundamentally cheaper because we share the output of DSA's lightweight indexer rather than full~$O(L^2)$ attention scores.
Second, we introduce systematic techniques for optimizing the sharing configuration, including a training-free greedy search to identify the optimal structural layout and a training-aware multi-layer distillation loss for parameter adaptation.
Although we instantiate IndexCache on DSA, the core principle extends to any sparse attention method that does not rely on a fixed sparse pattern but rather involves a dynamic token selection step: for instance, the block-level selection in MoBA~\citep{moba} and NSA~\citep{nsa} could similarly benefit from cross-layer reuse.